\documentclass[conference]{IEEEtran}
\usepackage[utf8]{inputenc}
\usepackage[T1]{fontenc}
\usepackage[spanish,es-tabla]{babel}
\spanishplainpercent
\usepackage{amsmath,amssymb,booktabs,url}
\usepackage[hidelinks]{hyperref}
\hypersetup{pdfauthor={Kevin Aguilar, Kenneth Campos y Kevin Cortés}}
\renewcommand{\IEEEkeywordsname}{Palabras clave}
\pagestyle{plain}
\raggedbottom
\newcommand{\pendiente}[1]{\par\medskip\noindent\textit{Pendiente de completar: #1.}\par\medskip}
\newcommand{\estado}{\noindent\textbf{Documento de trabajo. Validación experimental pendiente.}\par\medskip}
\newcommand{\autores}{\author{\IEEEauthorblockN{Kevin Aguilar, Kenneth Campos y Kevin Cortés}
\IEEEauthorblockA{Escuela de Ingeniería Electrónica\\Instituto Tecnológico de Costa Rica\\
EL4201 Procesamiento Electrónico de Potencia}}}
\newcommand{\referencias}{\begin{thebibliography}{9}\raggedright
\bibitem{enunciado} EL3313 y EL4201, \textit{Proyecto 4: instructivo y rúbricas}, secc. 4.3, 8.2 y anexo B, II semestre, 2026.
\bibitem{irf} Infineon, \textit{IRF540NPbF}, hoja de datos. [En línea]. \url{https://www.infineon.com/assets/row/public/documents/24/49/infineon-irf540n-datasheet-en.pdf}.
\bibitem{ir} Infineon, \textit{IR2110/IR2113: High and Low Side Driver}, hoja de datos. [En línea]. \url{https://www.infineon.com/assets/row/public/documents/24/49/infineon-ir2110-datasheet-en.pdf}.
\bibitem{acs} Allegro MicroSystems, \textit{ACS712}, hoja de datos, rev. 22. [En línea]. \url{https://www.allegromicro.com/-/media/Files/Datasheets/ACS712-Datasheet.ashx}.
\bibitem{pr} Imperix, \textit{Proportional resonant controller}, TN110, 2024. [En línea]. \url{https://imperix.com/doc/implementation/proportional-resonant-controller}.
\bibitem{plecs} Plexim, \textit{PLECS 5.0 User Manual}. [En línea]. \url{https://www.plexim.com/sites/default/files/plecsmanual-5-0.pdf}.
\bibitem{xadc} AMD, \textit{7 Series FPGAs and Zynq-7000 SoC XADC User Guide}, UG480, v1.11, 2022, cap. 2. [En línea]. \url{https://www.amd.com/content/dam/xilinx/support/documents/user_guides/ug480_7Series_XADC.pdf}.
\end{thebibliography}}

\hypersetup{pdftitle={Evaluación del inversor monofásico: documento de trabajo}}
\begin{document}
\title{Implementación y evaluación de un inversor monofásico con control PR de corriente}
\autores\maketitle\estado
\begin{abstract}
Se organiza la evaluación de un inversor monofásico con carga RL y control PR discreto. Se presentan los modelos, el método de medida y los resultados de simulación disponibles, junto con apartados identificados para parámetros reales, registros de laboratorio y comparación cuantitativa. La copia funcional incorpora tiempo muerto, resistencia de conducción y medición nominal cuantizada. Los resultados experimentales aún no están disponibles; no se declara validado el montaje ni la integración con MicroBlaze V.
\end{abstract}
\begin{IEEEkeywords}control de corriente, inversor, validación, PWM unipolar.\end{IEEEkeywords}

\section{Introducción}
El objetivo es seguir referencias sinusoidales de 50 Hz y 0,2, 0,5 y 1 A pico en una carga RL, con bus de 12 a 24 V \cite{enunciado}. El documento de diseño adjunto desarrolla la topología, el dimensionamiento y el candidato PR; este informe concentra implementación, resultados y discrepancias.

Se distinguen cálculos, simulaciones y mediciones. El diseño de Potencia corresponde a los tres autores; la ejecución con MicroBlaze V se coordina con Taller. Si se emplea una plataforma alternativa en la demostración propia, se identificará y se separarán sus resultados de los integrados.

\section{Montaje y parámetros}
La solución prevista utiliza cuatro IRF540N, dos IR2110 DIP-14 y un ACS712 de 5 A \cite{irf,ir,acs}. Dos resistencias de $3,3\ \Omega$, 5 W, forman la carga nominal. Una tercera permite pasar de 6,6 a 9,9 $\Omega$. El secundario del transformador se usa con el primario abierto y aislado.

\begin{table}[ht]\caption{Configuración y medidas por completar}\centering\small
\begin{tabular}{@{}lll@{}}\toprule
Magnitud & Modelo & Montaje\\\midrule
Bus & 12/24 V & Pendiente\\
R externa & $6,6\ \Omega$ & Pendiente\\
R de devanado & $0,6\ \Omega$ aprox. & Confirmar\\
L constante & 10 mH, supuesta & Verificar\\
PWM & 10 kHz & Pendiente\\
Control & 100 $\mu$s & Pendiente\\
Tiempo muerto & 1 $\mu$s & Pendiente\\
Límite de modulación & $\pm0,85$ & Confirmar\\\bottomrule
\end{tabular}\end{table}


Se midieron preliminarmente 9,291 mH a 120 Hz y 6,984 mH a 1 kHz. No se ha medido L a 50 Hz ni verificado su invariancia con la corriente.

\pendiente{esquemático completo y fotografía, auxiliares y protección instalados, identificación del devanado, equipos y límites de alimentación}

\section{Control y medición implementados}
El candidato PR usa $K_p=12$, $K_r=200$, resonancia en 50 Hz y $\omega_c=2\pi2$ rad/s. La estructura de ancho de banda finito y su discretización se sustentan en \cite{pr}. La ecuación conservada es
\begin{align}
r[k]={}&0,2509707453602995(e[k]-e[k-2])\nonumber\\
&+1,996504651666988r[k-1]\nonumber\\
&-0,9974902925463971r[k-2].
\end{align}
Se suma $12e[k]$, se normaliza por 12 V y se aplica una muestra después, limitando la modulación. La implementación final deberá confirmar período, adquisición y actualización efectiva.

La medición funcional usa 185 mV/A, offset de 2,5 V, divisor de 4,7 k$\Omega$/1 k$\Omega$, filtro de evaluación de 10 nF y cuantización de 12 bits \cite{acs,xadc}. Se supone calibración ideal, sin ruido ni deriva. El polo aproximado a 80 kHz no identifica el ancho de banda del módulo adquirido.

\pendiente{plataforma y versión, período y jitter reales, calibración, alineación ADC, retardos y tratamiento de saturación}

\section{Método experimental}
\subsection{Puesta en marcha}
Se documentarán instrumentos, sondas, referencias de tensión, escalas, ancho de banda y muestreo. Antes de energizar el puente se comprobarán las compuertas y el tiempo muerto. Se registrarán apagado, limitación de corriente y desconexión. El sensor se calibrará con una medida independiente.

\pendiente{capturas de compuerta, tiempo muerto medido, bootstrap, cero y sensibilidad del sensor e incertidumbre}

\subsection{Amplitud y fase}
Se ajustará, sobre ciclos completos,
\begin{equation}
i(t)\approx a\sin(\omega_0t)+b\cos(\omega_0t)+c.
\end{equation}
Así, $I_{1,pk}=\sqrt{a^2+b^2}$ y $\phi_i=\operatorname{atan2}(b,a)$. El desfase se expresará respecto de la fundamental de la referencia sincronizada. Con muestreo uniforme y ciclos completos puede utilizarse proyección seno/coseno; en otro caso se documentará el ajuste apropiado.
\begin{equation}
E_A=100\frac{|I_{1,pk}-I_{\mathrm{ref},pk}|}{I_{\mathrm{ref},pk}}.
\end{equation}
La ventana mínima es de cinco períodos completos de 50 Hz, equivalentes a 100 ms, después del establecimiento. No se sustituirá la fundamental por el máximo instantáneo ni por el RMS total.

\subsection{Rizado y transitorios}
El residuo de conmutación se medirá respecto de una componente de baja frecuencia estimada mediante un procedimiento documentado. Se indicarán filtro, banda, resolución y ventana local. Los límites nominales son $E_A\leq10\,\%$ y $\Delta i_{\mathrm{sw},pp}\leq0,20$ A \cite{enunciado}. Se reporta fase sin límite adicional.

Los escalones de amplitud y carga se realizarán por separado, conservando registros de referencia, corriente, tensión y PWM. Se analizarán picos, establecimiento y límites, distinguiendo transitorio y régimen.

\section{Simulación}
El modelo PR discreto conserva el puente ideal. La copia funcional añade Ron de 44 m$\Omega$ y órdenes con tiempo muerto de 1 $\mu$s \cite{irf,plecs}. No reproduce bootstrap, VGS, parásitos, temperatura, pérdidas de conmutación, caída o recuperación de diodos, saturación magnética ni anti-windup.
\input{datos_simulacion.tex}
La ventana nominal es [0,4; 0,5) s con exportación cada 2 $\mu$s. El máximo error nominal de la copia funcional fue aproximadamente 4,36\,\%. Los cambios de referencia se aplicaron a 0,205 s sin reiniciar la fase. La modulación alcanzó aproximadamente 0,836 en el cambio hacia 1 A, cercana al límite de 0,85.

Entre 0,100 y 0,101 s se comprobaron las órdenes cada 100 ns: hubo actividad en las cuatro señales, cero muestras de solapamiento y 40 huecos completos, con mínimo observado de 1 $\mu$s. Esto no verifica las compuertas físicas ni el PWM en RTL.

\section{Resultados experimentales}
Esta sección se completará con datos originales del montaje. Las celdas pendientes no representan valores nulos ni criterios aprobados.

\begin{table}[ht]\caption{Seguimiento experimental}\centering\small
\begin{tabular}{@{}cccccc@{}}\toprule
$V_{dc}$ & $I_{\mathrm{ref},pk}$ & $I_{1,pk}$ & $E_A$ & $\phi$ & $\Delta i_{pp}$\\
(V) & (A) & (A) & (\%) & ($^\circ$) & (A)\\\midrule
12 & 0,2 & Pend. & Pend. & Pend. & Pend.\\
12 & 0,5 & Pend. & Pend. & Pend. & Pend.\\
12 & 1,0 & Pend. & Pend. & Pend. & Pend.\\
24 & 0,2 & Pend. & Pend. & Pend. & Pend.\\
24 & 0,5 & Pend. & Pend. & Pend. & Pend.\\
24 & 1,0 & Pend. & Pend. & Pend. & Pend.\\\bottomrule
\end{tabular}\end{table}

\pendiente{registros de las tres referencias a 12 y 24 V, condiciones reales y ventanas de análisis}


\begin{table}[ht]\caption{Transitorios experimentales}\centering\small
\begin{tabular}{@{}lll@{}}\toprule
Ensayo & Pico de corriente & Establecimiento\\\midrule
0,2 a 0,5 A & Pendiente & Pendiente\\
0,5 a 1,0 A & Pendiente & Pendiente\\
R externa: $+50\,\%$ & Pendiente & Pendiente\\\bottomrule
\end{tabular}\end{table}

\pendiente{capturas, criterio de establecimiento, modulación y temperatura}

\section{Comparación y discusión}
Se compararán banco y PLECS en condiciones coincidentes. Una discrepancia de amplitud se puede expresar como
\begin{equation}
\delta_I=100\frac{I_{1,pk}^{\mathrm{exp}}-I_{1,pk}^{\mathrm{sim}}}{I_{1,pk}^{\mathrm{sim}}}.
\end{equation}
Se discutirán por separado calibración, ruido, cuantización, filtrado, retardo, tiempo muerto, semiconductores y variación de R y L. No se asignará una causa sin evidencia.

\pendiente{comparación cuantitativa, rizado, incertidumbre, pérdidas y explicación de discrepancias}

\section{Integración con Taller}
Se documentará la cadena ACS712--XADC--MicroBlaze V--PWM--IR2110 y los ensayos conjuntos. Esta integración corresponde a un rubro independiente de los 80 puntos propios de Potencia \cite{enunciado}.
\pendiente{interfaces y versiones finales, adquisición sincronizada y pruebas integradas; si no se completan, registrar preparativos e intentos}

\section{Conclusiones}
La simulación predice seguimiento nominal con error de amplitud inferior al 10\,\% bajo las hipótesis evaluadas, y recuperación después de los cambios de referencia y carga. La validación física y la comprobación final de rizado permanecen pendientes.
\pendiente{conclusiones basadas en mediciones, cumplimiento, limitaciones y mejoras}

\section*{Evidencias adjuntas}
Se acompañarán los modelos PLECS, scripts, parámetros, registros de simulación y datos experimentales originales, con instrucciones de reproducción. Cada figura deberá corresponder a un archivo y a condiciones identificadas. Los registros experimentales todavía no existen.
\newpage
\referencias
\end{document}
